\begin{korollar}{Normäquivalenzbasis im Endlichdimensionalen}
                {NormaequivalenzbasisImEndlichdimensionalen}
Seien $n \in \N$, $(E, \norm\cdot_E)$ ein $n$--dimensionaler normierter
$\K$--Vektorraum, $\norm\cdot$ eine Norm auf dem $\K^n$ und 
$d \coloneqq  d_{BM}((\K^n, \norm\cdot), (E, \norm\cdot_E))$.\\
Dann existiert eine Basis $x \in E^n$ mit der Eigenschaft
\[\frac 1 d \norm\cdot \le\norm\cdot_E \circ \bprod\cdot x_{n, E} \le \norm\cdot\text.\]
\end{korollar}
\begin{proof}
Da $E$ und $\K^n$ endlichdimensional sind und die gleiche Dimension haben,
existiert nach~\ref{ExistenzEinerBanachMazurAbbildung} eine
Banach--Mazur--Abbildung 
$T \in \Inv(\K^n, E)$ zwischen $(\K^n, \norm\cdot)$ und $(E, \norm\cdot_E)$,
also sind 
\begin{align}\opnorm T{\K^n}E=1 \quad \text{ und }\quad 
\opnorm{T \inv}E{\K^n} = d\text. \label{nieeq:1}\end{align}
Seien $e \coloneqq  (e_i^n)_{i=1}^n$ und $x \coloneqq  \pfeil T n (e)$. Offenbar ist $x$ eine Basis von $E$.\\
Sei $t \in \K^n$. Dann gilt:
\[ \bprod t x _{n, E}= \bprod t {\pfeil T n (e)} _{n, E}
\overset{\refclaim{EigenschaftenDesLinearkombinators}
{EigenschaftenDesLinearkombinators2}}{=} 
T\left(\bprod t e _{n, \K^n}\right) = T(t)\text.\]
Somit ist also $T = \bprod \cdot x_{n, E}$ und daher 
\begin{align}\norm\cdot_T = \norma\cdot_E \circ \bprod \cdot x _{n, E}\text.
\label{nieeq:2}\end{align}
Wegen der endlichen Dimension sind $(E, \norm\cdot_E)$ und $(\K^n, \norm\cdot)$
vollständig, $T$ ist injektiv und hat abgeschlossenes Bild.\\
Daher liefert~\ref{EigenschaftenDerBildhalbnorm}:
\begin{align}\inf\menge{ \norm{Ta}_E }{a \in \kran {\K^n} 0 1} \cdot \norm\cdot
    \le \norm\cdot_T \le \opnorm TE{\K^n}\cdot \norm\cdot\text. \label{nieeq:3}
    \end{align}
Auch die Voraussetzungen von~\ref{OperatornormDerInversen} sind erfüllt,
sodass wir weiterhin erhalten:
\begin{align}\label{nieeq:4}\inf\menge{ \norm{Ta}_E }{a \in \kran {\K^n} 0 1} 
 = \frac 1 {\opnorm{T\inv}E {\K^n}} \overset{\eqref{nieeq:1}}= \frac 1 d\text. \end{align}
Insgesamt folgt:
\begin{align*}
\frac 1 d \norm\cdot &\overset{\eqref{nieeq:4}}=
 \left(\inf\menge{\norm{Ta}_E }{a\in\kran{\K^n}01}\right) \cdot \norm\cdot \\
 &\overset{\eqref{nieeq:3}}\le \norm\cdot_T \overset{\eqref{nieeq:2}}= 
 \norm\cdot_E \circ \bprod \cdot x _{n, E}
 \overset{\substack{\eqref{nieeq:2},\\\eqref{nieeq:3}}}\le 
 \opnorm T {\K^n} E \norm\cdot \overset{\eqref{nieeq:1}}
 = \norm \cdot \text.
\end{align*}
\end{proof}